\chapter{Anforderungen}
\label{chap:anforderungen}


Die Anforderungen dieser Arbeit orientieren sich am \gls{gls-ETSI}-\gls{gls-NFV}-Framework \parencite{ETSI-NFV} und an Standards für Zuverlässigkeit und Verfügbarkeit virtualisierter Netzwerk-Funktionen \parencite{ETSI-REL}. Im Folgenden wird zwischen der prototypischen Implementierung und den in der Literatur empfohlenen produktiven Vorgaben unterschieden.

\section{Funktionale Anforderungen}

\subsubsection{Prototypische Implementierung}

Der Prototyp basiert auf den folgenden funktionalen Anforderungen für eine containerbasierte Architektur und \gls{gls-KI}-basierte Sprachinteraktion:

\begin{itemize}
  \item \textbf{\gls{gls-SIP}-Gateway:} RFC 3261-konformer Rufaufbau und Call-Control über ein lokales Asterisk-\gls{gls-PBX} \parencite{rfc3261}
  \item \textbf{Media-Bridge:} Transcoding zwischen Audio-Codecs (G.711 für \gls{gls-PSTN}, Opus für \gls{gls-WebRTC}, \gls{gls-PCM16} für \gls{gls-ASR}) und asymmetrisches \gls{gls-RTP}-Handling für \gls{gls-SIP} $\leftrightarrow$ \gls{gls-WebRTC} \parencite{rfc3550}
  \item \textbf{\gls{gls-KI}-Pipeline:} \gls{gls-ASR} mit \gls{gls-VAD} für Spracherkennung, \gls{gls-LLM} für Intent-Erkennung und Dialogverwaltung, \gls{gls-TTS} für Antwort-Generierung
  \item \textbf{Sitzungsverwaltung:} Statusverfolgung für aktive Calls mit minimaler Persistenz
\end{itemize}

Der detaillierte Signalfluss wird in Kapitel \ref{chap:architektur} beschrieben.

\subsubsection{Produktive Vorgaben nach Standards}

Produktive Systeme erweitern die Basisfunktionalität um mehrere empfehlenswerte Anforderungen. Basierend auf dem \gls{gls-ETSI}-\gls{gls-NFV}-Rahmenwerk \parencite{ETSI-NFV, ETSI-REL} und bewährten Praktiken für containerbasierte Asterisk-Deployments empfiehlt sich eine produktive Architektur mit folgenden Vorgaben \parencite{Antonova2025}:

\begin{itemize}
  \item \textbf{Verteilte Systemkomponenten:} Trennung von \gls{gls-SIP}-Signalisierung (Kamailio), Medienverarbeitung (Asterisk Core), Datenspeicherung (PostgreSQL) und Lastverteilung als separate, unabhängig skalierbare Container-Services \parencite{Antonova2025}
  \item \textbf{Fehlertoleranz und Verfügbarkeit:} Automatische Neu-Initialisierung fehlerhafter Container, automatische Wiederherstellung von Diensten, internes Cluster-Load-Balancing \parencite{Antonova2025}
  \item \textbf{Horizontale Skalierbarkeit:} Automatische Pod-Erzeugung basierend auf \gls{gls-CPU}- und \gls{gls-RAM}-Metriken mittels Horizontal Pod Autoscaler, verteilte Lastverteilung zwischen Asterisk-Replikas \parencite{Antonova2025,KubernetesDocs}
  \item \textbf{Persistierung und Verwaltung:} Zentralisierte Datenspeicherung (\gls{gls-CDR}, Konfiguration, Authentifizierung) in relationalen Datenbanken, deklarative Konfigurationsverwaltung über Kubernetes-Manifeste \parencite{Antonova2025,KubernetesDocs}
  \item \textbf{Observability und Automatisierung:} Monitoring von System-Metriken und Logs, automatisierte Deployment- und Rollback-Prozesse, DevOps-Integration für kontinuierliche Bereitstellung \parencite{Antonova2025}
\end{itemize}

\section{Nicht-funktionale Anforderungen}

Der Prototyp basiert auf den folgenden nicht-funktionalen Anforderungen für eine containerbasierte Architektur und \gls{gls-KI}-basierte Sprachinteraktion:

\subsubsection{Prototypische Implementierung}

\begin{itemize}
  \item \textbf{Latenz:} End-to-End 2--5\,s für \gls{gls-KI}-Dialog-Interaktionen (akzeptabel für asynchrone Konversation, nicht \gls{gls-ITU-T}-konform für klassische Telefonie \parencite{ITU-T-G114}).
  \item \textbf{Verfügbarkeit:} Best-Effort ohne garantierte Uptime, kein Failover, Single-Host-Deployment.
  \item \textbf{Skalierbarkeit:} Vertikales Scaling begrenzt auf Host-Ressourcen, kein horizontales Scaling.
  \item \textbf{Sicherheit:} Container-Isolation via Docker, Basis-Authentifizierung für Zugang zu Verwaltungsschnittstellen, keine Verschlüsselung für Entwicklungs-/Testzwecke.
  \item \textbf{Datenschutz:} Minimale Session-Persistenz, keine langfristige Datenspeicherung.
\end{itemize}

\subsubsection{Produktive Vorgaben nach Standards}

Die folgenden Anforderungen beschreiben das produktive Zielmodell nach \gls{gls-ETSI}-REL-Standard \parencite{ETSI-REL} und \gls{gls-ITU-T} G.114 \parencite{ITU-T-G114}. Sie folgen bewährten Praktiken anhand containerbasierten Asterisk-Deployments \parencite{Antonova2025}:

\begin{itemize}
  \item \textbf{Latenz:} < 400\,ms für Media-Übertragung (Sprachübertragung), < 2\,s für \gls{gls-KI}-Dialog-Systeme (asynchrone Interaktion), nach \gls{gls-ITU-T} G.114 \parencite{ITU-T-G114}.
  \item \textbf{Verfügbarkeit:} Hohe Verfügbarkeit mit Redundanz kritischer Komponenten und automatischem Failover \parencite{ETSI-REL}.
  \item \textbf{Skalierbarkeit:} Horizontale Skalierung über Kubernetes mit automatischer Last-Verteilung basierend auf Standard-Metriken wie \gls{gls-CPU} und \gls{gls-RAM} oder benutzerdefinierten Metriken wie Anrufvolumen \parencite{KubernetesDocs}. Containerbasierte Telefonsysteme können solche \gls{gls-VoIP}-spezifischen Metriken über benutzerdefinierte Monitoring-Lösungen realisieren \parencite{Antonova2025}.
  \item \textbf{Sicherheit:} Verschlüsselte Datenübertragung für kritische Komponenten.\gls{gls-SRTP} für Medien anhand von RFC 3711 \parencite{rfc3711}, Schlüssel-Management für \gls{gls-API}-Schlüssel und sichere Aufbewahrung von Zugangsdaten außerhalb des Quell-Repositories \parencite{Antonova2025}.
  \item \textbf{Datenschutz:} \gls{gls-DSGVO}-konforme Praktiken einschlie\ss lich Datenminimierung, Pseudonymisierung und dokumentierter Aufbewahrungsrichtlinien \parencite{Data2025_AIDataQuality}.
\end{itemize}


